% Expected vs Actual Behavior
\needspace{10cm}
\section{Expected vs Actual Behavior}

% Clearly distinguish what should happen from what does happen.
% Include concrete, observable outcomes — not interpretations.

\subsection{Comparison}

\begin{wyrdtab}{l p{5.5cm} p{5.5cm}}
\textbf{Aspect}     & \textbf{Expected}                     & \textbf{Actual} \\
Behavior            & Describe correct behavior              & Describe observed behavior \\
Output / Response   & Expected output or status code         & Actual output or error \\
Data State          & Expected database or state change      & Actual data state \\
Performance         & Expected response time or throughput   & Actual measured values \\
\end{wyrdtab}

\subsection{Error Output}

% Include the exact error message, stack trace, or log output.
% Use the code environment to preserve formatting.

\begin{criticalpoint}
Paste the primary error message or exception here for quick reference.
\end{criticalpoint}

% Uncomment and populate with the full stack trace or log excerpt:
%
% \begin{wyrdcode}[Stack Trace]{text}
% Traceback (most recent call last):
%   File "module.py", line 42, in function_name
%     result = operation()
% TypeError: description of the error
% \end{wyrdcode}

\subsection{Visual Evidence}

% Include screenshots or recordings if the defect has a visual component.
% Place image files alongside the project and reference them here.

% \begin{figure}[h]
%     \centering
%     \includegraphics[width=0.8\textwidth]{sections/screenshot.png}
%     \caption{Description of what the screenshot shows}
% \end{figure}
